\documentclass[instructions.tex]{subfiles}
\begin{document}

\chapter{Combien le penchant qu'on a au péché est à craindre.}


Tous les hommes ont du penchant au mal : les uns, à la volupté ; les autres, à l'intérêt ; d'autres, à la colère, à la jalousie, à la vengeance ; plusieurs, à la gourmandise et à l'oisiveté ; presque tous, à l'indépendance, à l'orgueil et à la gloire. Si nous n'avions que le penchant au mal, nous ne serions pas encore coupables, puisque ce penchant n'est pas un péché, mais une suite de notre malheureuse origine. Les saints ont éprouvé ce penchant : il a été pour eux, et il doit être pour nous un sujet de gémissement, une occasion de combats et de mérites. Nous ne sommes donc pas criminels de ressentir nos mauvais penchans ; mais nous sommes coupables de ne vouloir pas les connaître et de les suivre.

Connaître ses penchans déréglés et les combattre, voilà ce que saint Paul appelle vivre selon l'esprit ; c'est en vivant de la sorte qu'on gagne le ciel. Suivre ses mauvais penchans et négliger de les connaître, c'est ce qu'on appelle vivre selon la chair ; vivre de la sorte, c'est s'égarer et se perdre. Doit-on s'étonner de voir tant de personnes vivre dans l'égarement, puisqu'il y en a si peu qui s'appliquent à connaître les penchans qui sont le mobile de leur conduite ? connaissance néanmoins nécessaire pour le salut. En suivant ses penchans vicieux, on fait des fautes plus fréquentes, plus grièves, plus irréparables.


\primo On pèche plus fréquemment, parce que le penchant nous fait penser souvent aux objets que nous aimons ou que nous haïssons. D'où viennent tant de pensées impures, si ce n'est de l'inclination que vous avez aux sales voluptés ? D'où viennent ces projets qui roulent sans cesse dans votre esprit sur les richesses, sur les moyens d'acquérir, si ce n'est de votre avarice, de votre penchant à l'intérêt ? D'où viennent tant de pensées contre la charité, de murmures, de mépris des autres, si ce n'est de votre antipathie pour certaines personnes, de votre aversion à souffrir, et du penchant que vous avez à l'orgueil ? D'où viennent ces fréquentes dissipations dans vos prières, si ce n'est de vos penchans, qui font naître tantôt des pensées flatteuses sur les objets que vous aimez, tantôt des pensées importunes sur les objets que vous avez en aversion ?

Des pensées suivent les autres pensées, les désirs, les paroles, les saillies d'humeur, les actions, parce que nos mauvais penchans, que saint Jacques appelle la concupiscence, en inspirant les pensées, font naître le péché ; le péché, étant consommé dans le cœur par la délibération, donne la mort\footnote{Concupiscentia, cum conceperit, parit peccatum ; peccatum verò, cum consummatum fuerit, generat mortem. Jac. 1.}. Il est donc important de connaître et de réprimer nos penchans déréglés, puisqu'ils nous portent au mal presque à tout moment.


\secundo Une seconde raison, c'est qu'en les suivant on pèche pour l'ordinaire plus grièvement. Il est bien dangereux, lorsqu'on ne pense qu'à se contenter et à satisfaire son inclination, qu'on ne s'expose ou qu'on n'aille jusqu'au péché mortel. Tout paraît permis à celui qui juge des choses par passion. Judas se croyait peut-être innocent, et bien d'autres jugeraient de même, lorsqu'il disait qu'il valait mieux donner aux pauvres le prix d'un baume, que de le répandre sur les pieds du Sauveur. Mais l'Évangile dit qu'il parlait de la sorte, parce qu'il était avare et qu'il aimait l'argent plus que les autres. Judas, comme beaucoup d'autres aujourd'hui, ne prenait pas garde que ce penchant à l'intérêt souillait son cœur par l'avarice. Aurait-il cru lui-même que ce penchant le porterait au plus exécrable des crimes, à vendre son maître ?

Que si la conscience nous fait connaître qu'il y a du mal en certaines choses, le penchant qui aveugle, fait croire qu'il n'y en a pas beaucoup. Mais l'aveuglement et l'erreur qui viennent de la passion ou du penchant, excusent-ils devant Dieu ? Nous croyons n'être pas coupables, nous nous excusons en disant : Je n'y pensais pas, je ne savais pas ; mais Jésus-Christ, qui nous commande de veiller sur nous, d'examiner les mouvemens de notre cœur, nous excuse-t-il de même ?

Combien de gens favorisent en tout leur penchant pour leur intérêt et leurs plaisirs, qui ne consultent point sur des choses essentielles, qui s'appuient sur des décisions extorquées pour justifier leur conduite, leurs injustices, leurs scandales ; qui se croient innocens, tandis que Dieu les condamne ! Pourquoi cherchez-vous à justifier vos égaremens, disait le Seigneur à son peuple ? Vous avez dit : Je suis innocent, je n'ai point fait de mal. Je vais entrer en jugement avec vous, parce que vous avez dit : Je n'ai point fait de mal\footnote{Dixisti : Absque peccato et innocens ego sum. Ecce ego judicio contendam tecum, eo quod dixeris : Non peccavi. Jer. 2.}.


\section{Résolutions}

\primo Je m'appliquerai sérieusement à connaître mes mauvais penchans.

\secundo Je leur ferai une guerre continuelle, en résistant aux impulsions qu'ils me donneront, dès que je m'en apercevrai.

\tertio Et je demanderai souvent à Dieu la grâce de les combattre efficacement.

Mon Dieu, faites-moi la grâce de connaître les inclinations au mal qui dominent en moi, et de les réduire entièrement sous le joug de votre loi sainte.

\end{document}
